% ==========================================================================
% GENERATED by scripts/split_paper.py. Do not edit this file.
%
% Source: imported_papers/Neural_Computation__CNA/NECO-D-24-00304-Source.tex
% Part 5: Relation to Alternative Concepts
%
% Edits belong upstream, in the published paper. To change what the
% chapter includes, change imports/cna.toml and re-run
% the splitter.
% ==========================================================================
\section{Relation to Alternative Concepts} \label{sec:RelatedWork} 
The central issue behind our model is how the structure of the environment is represented. In biology, this structure would be accessed by the senses in the form of sensory patterns --- large active subsets of sensory elements (e.g., retinal ganglion cells). While global sensory patterns vary across different conditions and viewpoints, local sensory structure tends to exhibit consistency. For instance, the sensory pattern associated with a static object may appear different under varying eye positions or perspectives, yet local fragments of the pattern remain stable and coherent (we demonstrate this in our experiments by having constant straight-line patterns that are stabilized, while changing Gaussian noise is ignored).
The task of learning -- both in biology and in a model like ours -- is to pick up durable structures from sensory patterns in spite of their fleeting nature. We proceed here with a simple assumption regarding this durable structure: There is an important class of coherent sensory patterns that are composed of a finite set of overlapping pattern fragments. We assume that these pattern fragments can be represented by the recurrent connectivity of net fragments and that this connectivity can be learned from the statistics of sensory patterns with the help of Hebbian plasticity.  This is the background against which we discuss some models of neural computation currently dominating the literature.


\subsection{Associative Memory} \label{sec:associative_memory}

The central concept in the classic works of \cite{hebb_organization_1949}  and \cite{hayek} is the assembly, a set of neurons supporting each other's activity by mutual excitatory connections, connections that have been created by synaptic, ``Hebbian'', plasticity.  This concept has found a concrete mathematical formulation as ``associative memory'' (AM) \citep{amari_learning_1972, Hopfield1982}, where stored patterns are encoded by enhancing excitatory connections between all pairs of co-active neurons.

While AM mechanisms are effective for storing and retrieving known patterns, they often fail when applied to perceptual tasks such as object recognition, where previously unseen structures must be classified based on prior experience.
This limitation arises because AM tends to store complete patterns as static ``snapshots.'' 
The reason for this rigidity of AM states is their mode of recording: The excitatory connections between all pairs of neurons active at a particular moment are strengthened simultaneously, without regard to whether a connection is significant in the sense of applying to many patterns (that is, being part of a repeating pattern fragment) or insignificant, applying only to this particular pattern.
  
In our model, in contrast, the synaptic weights are gradually modified in response to the statistical properties of a large ensemble of sensory patterns. According to our Eq. \ref{eq:hebbian_plasticity}, a particular synapse can only grow if the pair of neurons it connects is simultaneously active again and again (indicated by positive correlation $\boldsymbol{\rho}_{avg}$), or it decays due to negative correlation.
This processing allows the system to gradually extract and store re-occurring (statistically relevant) sub-patterns rather than memorizing isolated global patterns.

A functional restriction on AM is that different stored states have to be statistically independent, that is, as sets of active neurons, they have to have low pairwise overlap. This is an awkward constraint, as sensory patterns are composed of repeating pattern fragments and thus are bound to have a large overlap. In classical AM, this constraint has been addressed just by requiring patterns to be small subsets of all neurons \citep{Tsodyks2007, Palm2013}. Our model can address the small-overlap constraint in a systematic way with the help of multiple neurons (``competing neurons'') that respond to the same trigger feature but compete with each other for activity so that pairs of pattern fragments that have high feature-wise overlap can select different competing neurons, resulting in the corresponding net fragments having low neuron-wise overlap. This is made possible by the statistical mode of strengthening synapses and is not possible in AM with its one-shot learning mode.

Our model also addresses an open issue of AM: The selection of coherent sensory patterns for storage. As coherence is defined as being composed of a fixed set of overlapping pattern fragments (stored in net fragments), it is only natural to require that the learning process converges to a connectivity state in which only sets of neurons that segment the input into coherent sensory patterns can be activated stably \citep{Malsburg2024}.

The canonical formulation \citep{amari_learning_1972,Hopfield1982} of the AM's state dynamic defines an ``energy'' that acts as a Lyapunov function. Beginning from any initial state, neurons are visited sequentially, and their state is flipped to gradually reduce the energy (simultaneous switching of the neurons would lead to non-sensical oscillations). Our model's dynamics is not based on a Lyapunov function and a time-consuming gradual process, relying instead on an externally imposed periodic attenuation control parameter that lets the set of neurons originally activated by the sensory input collapse to a subset formed by neurons that support each other by excitatory connections.

\subsection{Boltzmann Machine} \label{boltzmann_machine}

Regarding Boltzmann Machines (BM), we refer to \citet{Hinton1983, hinton_constraint1984}, taking the latter as the authoritative reference. Both the BM and our model describe networks of recurrent connections between neurons, where some neurons get activated by sensory input (``visible units''), and by employing synaptic plasticity, both approaches aim to turn the system into a constraint satisfaction network that uses inner connections to stabilize the patterns suggested by the sensory input.
  
The BM as described in \citet{hinton_constraint1984} and our model differ, however, in a number of ways.  Inspired by statistical mechanics, the BM speaks of equilibrium macro-states, each of which is realized as an ensemble of micro-states, the probability of a micro-state being proportional to a negative exponential of a state ``energy'' (Boltzmann distribution), which is defined as in Amari-Hopfield's Associative Memory as a function of neural activities and connection weights. The BM further assumes that the system compares two modes, clamped (in which the sensory neurons are determined by the input) and free running (no sensory input), and lets synaptic plasticity equalize the signal correlations of pairs of neurons observed in one of the two modes with those observed in the other.
In our model, neural activity is not modeled in terms of probabilistic ensembles of states but
as discrete-time trajectories over a fixed number of time steps $T$, during which the same sensory input is constantly refined.
In distinction to the Boltzmann Machine, these trajectories evolve deterministically, according to update rules that depend on both the input and recurrent connectivity.
Importantly, the dynamics are externally modulated: At each step, a control mechanism attenuates the activity of neurons that lack sufficient recurrent support from other active neurons;
see Eqs. \ref{eq:inhibition} and \ref{eq:s2_final_output}.
In contrast, Boltzmann Machines rely on stochastic sampling over ensembles of network states, governed by a distribution derived from an energy function. While the statistical mechanics framework underlying Boltzmann Machines is elegant and conceptually rich, it is computationally inefficient in settings where a large number of intermediate states must be generated and compared to approximate equilibrium. Our model avoids this bottleneck by converging rapidly through deterministic dynamics over a small fixed number of steps.

A necessary precondition for the feasibility of the goals of the BM and our model is the existence of a system of constraints on the sensory patterns that can be learned from a finite number of such patterns and which then suffice to stabilize all future patterns.  Given the bias-variance dilemma \citep{Geman1992} and the no-free-lunch theorems \citep{NoFreeLunch}, a learning system's architecture needs to drastically constrain its search space by a bias that is tuned to the application domain \citep{learning1997tom}. No such bias is built into the BM's architecture in its generality \footnote{The concrete application described in \citet{hinton_constraint1984} (the `Encoder Problem') provides such bias, but is trivial in the sense of explicitly defining the full (and very small) set of sensory patterns.}, while our model's architecture explicitly incorporates the bias that the sensory patterns of its application domain, vision, are all \emph{coherent composites of a finite set of pattern fragments}, a bias that is effective yet captures an essentially infinite number of sensory patterns.

The realization that the BM in its generality would imply unfeasible learning times led to the introduction of Restricted Boltzmann Machines (RBM) \citep{Smolensky1987,Hinton2006}. RBMs consist of two layers, one visible and one latent, with bidirectional connections between them, but no connections within them. A stack of RBMs is used in \citep{Hinton2006}
to construct an autoencoder in which RBM learning is used to pre-train the connections between layers in a sequence from bottom (input) to top (the sparse coding layer). Then, the whole encoder-decoder stack is trained end-to-end by backpropagation of error.
We discuss the differences between our model and networks trained by backpropagation of error (including RBMs) in the next sub-section.

\subsection{Artificial Neural Networks (ANN)} \label{sec:ANN}

Artificial neural networks (ANNs), as concerned here, are layered structures with feedforward layer-to-layer connections from input to output, and these connections are trained by back-propagation of error \citep[p 292]{rosenblatt_principles_1962}. Prominent examples are multilayer perceptrons (MLPs) \citep{prince2023understanding}, CNNs \citep{lecun1998gradient}, autoencoders \citep{Hinton2006}, and transformers \citep{attention2017}.  In these systems, individual neurons in one layer represent
%significant pattern fragments
feature compositions % I think this is more precise in this context?
that occur in the layer below.  The complexity and variability of features represented by individual neurons grow from layer to layer  (as shown in reconstructions, see \citet{Zeiler2013}). The hierarchy of layers thus constitutes a hierarchy of nested pattern fragments. 

The architectures of both ANNs and our model implement the bias that a virtually infinite set of potential input patterns can be represented in the form of flexible compositions of a finite set of pattern fragments.  They differ fundamentally, however, in how these pattern fragments are implemented, how they are combined, and how they are learned. 

Beginning with the latter, ANNs select the features to be represented on the basis of their effect on lowering the system's prediction error, whereas in our model the selection of features to be represented by net fragments is entirely a matter of the statistics of input patterns (as is, incidentally, the case with the RBM pretraining of layer connections in \citet{Hinton2006}). 

Regarding implementation, ANNs represent elementary pattern fragments in the first latent layer by feature neurons. In our model, which has only one layer, elementary pattern fragments are represented by minimal net fragments, each composed of a neuron and a number of excitatory connections from other neurons within the same layer. That a neuron remains active under attenuation indicates that it has recognized the pattern formed by these other also activated neurons. ANNs represent a nested hierarchy of pattern fragments, as higher-level neurons respond to combinations of active lower-level neurons.  Our model equally represents a nested hierarchy of pattern fragments within one layer, in this case composed of a neuron, its immediate supporters, the supporters of the supporters, and so on. 

The rules for composing larger structures out of smaller pattern fragments correspondingly differ between ANNs and our model. In CNNs, this flexibility is mostly implemented by max pooling, or in transformers by the attention mechanism that selects active neurons in one layer to contribute to the activity of neurons in the next. This freedom of choosing pattern fragments in lower levels is exerted independently by the neurons in a given level, without regard to the mutual compatibility of those selected lower-level patterns. This lack of coherence is even more drastic between ANNs' high-level feature neurons, which don't express in any way mutual overlap in terms of constituent sub-patterns (the binding problem, see \citet{Greff2020}), the likely cause of vulnerability of ANNs to adversarial attack \citep{Leadholm2022}. In our model, in contrast,  the net fragments that respond to different pattern fragments have to fit together into one consistent net, in which each active neuron is sufficiently supported by other active neurons. This expresses the bias that significant patterns in the visual environment are composed of pattern fragments {\it in a coherent way}. We consider this as the solution to the binding problem \citep{Greff2020} and as an explanation of the Gestalt phenomenon \citep{Wagemans2012}.
